\section{Conclusions}\label{sec:conclusion}
We have constructed a single-metric relativistic modified-gravity cosmology combining an AeST scalar--vector sector with an auxiliary Type-II acceleration sector, while fixing particle cold dark matter and the bare cosmological constant to zero. The scalar kinetic function contains a current-linked nonseparable operator that leaves both the homogeneous scalar background and screened static galactic function unchanged but makes the radiation-era scalar system strictly damped on the physical branch. The same construction has a closed perturbation map in CLASS, a consistent nonlinear constraint count, a canonically controlled cubic scalar sector, and a regular first-order Dirac reduction.

After the theory source was fixed, the final Planck 2018 + DESI DR2 + Pantheon+ likelihood evaluation gave $\chi^2=2422.69597$, lower by $6.2232$ than the matched optimized $\LCDM$ reference. The corresponding AIC difference favors the modified-gravity realization, whereas the adopted BIC convention slightly favors the lower-dimensional reference. An independent particle-CDM profile places the exact zero-CDM realization within $\Delta\chi^2<1$ of its sampled minimum. The result therefore establishes a stable and observationally competitive zero-particle-CDM realization, while stopping short of claiming that current data uniquely require the absence of particle dark matter.

The source code, exact parameter file, numerical summaries, figure-generation scripts, and reproducibility hashes are prepared for public release so that each principal result can be independently retested.
